% 原试卷 PDF 第46～50页。
\begin{question}[index=17,points=10]
求曲线 $\displaystyle y=\frac{1+\displaystyle\int_0^{\frac1x}\mathrm e^{t^2}\dd t}
{1-\mathrm e^{\frac1x}}\ (x>0)$ 的斜渐近线。
\end{question}
\QuestionBreak\WorkSpace

\begin{question}[index=18,points=12]
设 $y=y(x)\ (0\leq x\leq1)$ 具有二阶连续导数，曲线 $y=y(x)$ 为凸的，且位于第一象限，在其上每一点的曲率等于
$\displaystyle\frac{y}{(2-y^2)^{\frac32}}$，又设该曲线在原点与直线 $y=x$ 相切，求曲线的表达式。
\end{question}
\QuestionBreak\WorkSpace\WorkPageBreak

\begin{question}[index=19,points=12]
已知连续函数 $f(x)$ 满足 $\displaystyle f(x)=\mathrm e^x-\int_0^x f(t)\dd t$，将 $f(x)$ 展开成 $x$ 的幂级数。
\end{question}
\QuestionBreak\WorkSpace

\begin{question}[index=20,points=12]
计算曲面积分 $\displaystyle\iint_\Sigma x^2\dd y\dd z+\dd z\dd x+
\frac{\mathrm e^2}{\sqrt{x^2+y^2}}\dd x\dd y$，其中 $\Sigma$ 为 $z=\sqrt{x^2+y^2}\ (0\leq z\leq1)$ 的下侧。
\end{question}
\QuestionBreak\WorkSpace\WorkPageBreak

\begin{question}[index=21,points=12]
若矩阵 $\symbfit A$ 的伴随矩阵
$\symbfit A^*=\begin{pmatrix}3&2&-2\\0&-1&0\\a&2&-3\end{pmatrix}$
相似于矩阵 $\symbfit B=\begin{pmatrix}-1&0&0\\-2&1&0\\0&0&-1\end{pmatrix}$，其中 $|\symbfit A|>0$。
\begin{enumerate}[label=（\arabic*）,leftmargin=2.5em,labelsep=.4em,itemsep=3pt,topsep=3pt]
\item 求 $a$ 的值；
\item 求 $\symbfit A^{99}$．
\end{enumerate}
\end{question}
\QuestionBreak\WorkSpace

\begin{question}[index=22,points=12]
设商场某日的顾客总数 $N$ 服从泊松分布 $P(\lambda)$，单个顾客每次的消费额为 $1$（单位：元），令 $X_j$ 表示第 $j$ 个顾客 $n$ 次试验消费的次数，$X_j$ 与 $N$ 独立且服从二项分布 $B(n,p)$．
\begin{enumerate}[label=（\arabic*）,leftmargin=2.5em,labelsep=.4em,itemsep=3pt,topsep=3pt]
\item 已知 $N=m\ (m>0)$ 的条件下，求全天营业额 $S$ 的概率分布；
\item 计算 $E(S\mid N)$；
\item 计算 $ES$．
\end{enumerate}
\end{question}
\WorkSpace
